% J34 — F_p Extensions of CL_BHML: Universality Across Six Prime Fields
%
% Brayden R. Sanders, M. Gish, 7Site LLC
% Submitted to: Communications in Algebra
% Source corpus: Variant Catalog "TSML F_p extensions" (Atlas/LENS_TAXONOMY_2026-05-06/VARIANT_CATALOG.md)
%                + WP118_FP_UNIVERSALITY.md (sprint18_bridge_dirac_2026_05_04)
%                + GAP_AUDIT.md
% Companions cited as already-submitted: J21
% MSC 2020: 17A30, 11T55, 17A40, 12E20
% Tier: B (extension of J21's TSML F_p universality to the BHML substrate)

\documentclass[11pt,reqno]{amsart}

\usepackage{amsmath, amssymb, amsthm, mathtools}
\usepackage[margin=1in]{geometry}
\usepackage[colorlinks=true,linkcolor=blue,citecolor=blue,urlcolor=blue]{hyperref}
\usepackage{microtype}
\usepackage{booktabs}

\theoremstyle{plain}
\newtheorem{theorem}{Theorem}[section]
\newtheorem{proposition}[theorem]{Proposition}
\newtheorem{lemma}[theorem]{Lemma}
\newtheorem{corollary}[theorem]{Corollary}

\theoremstyle{definition}
\newtheorem{definition}[theorem]{Definition}
\newtheorem{conjecture}[theorem]{Conjecture}

\theoremstyle{remark}
\newtheorem*{remark}{Remark}

\newcommand{\Z}{\mathbb{Z}}
\newcommand{\F}{\mathbb{F}}
\newcommand{\TSML}{\mathrm{TSML}}
\newcommand{\BHML}{\mathrm{BHML}}
\newcommand{\Aut}{\mathrm{Aut}}
\DeclareMathOperator{\HARM}{H}
\DeclareMathOperator{\VOID}{V}

\title[F_p Extensions of CL_BHML]
{F_p Extensions of CL\_BHML: Universality Across Six Prime Fields}

\author{Brayden R.~Sanders}
\address{7Site LLC, Hot Springs, Arkansas, USA}
\email{brayden@7site.co}

\author{M.~Gish}
\address{Independent Researcher, Hot Springs, Arkansas, USA}
\email{monica.gish1992@gmail.com}

\subjclass[2020]{17A30, 11T55, 17A40, 12E20}
\keywords{commutative non-associative algebra, BHML composition table,
finite-field extension, prime universality, idempotent count,
$\Z/10\Z$ substrate}

\date{\today}

\begin{document}

\begin{abstract}
The $\mathrm{CL}_\BHML$ ($28$-HARMONY) composition lattice on
$\Z/10\Z$ is the second canonical substrate in the three-substrate
architecture (alongside $\mathrm{CL}_\TSML$ and $\mathrm{CL}_\mathrm{STD}$).
Its restriction to the $4$-core $\{0,7,8,9\}$ admits an extension to
the $\F_p$-bilinear $4$-dimensional algebra $V_p$ for any prime $p$.
We extend the F_p universality result of Sanders--Gish
\cite{SandersGishFpTSML} (J21, the $\TSML$ side) to the parallel
$\BHML$ side, verifying that the structural skeleton of $V_p^\BHML$
(idempotent count, automorphism orbit structure, eigenspace
signatures of left-multiplication, power-associativity, associator
image localization) is the \emph{same across all six primes}
$p \in \{2, 3, 5, 7, 11, 13\}$ tested.

The $\BHML$-side BHML\_8\_YM determinant identity
$\det(\BHML_8^{\text{YM}}) = +70 = \binom{8}{4}$ holds over $\Z$ and
is preserved (in the $\F_p$ versions where it makes sense) across the
verified primes. The two parallel chains (TSML\_$k$ and BHML\_$k$ for
$k \in \{1, 4, 5, 6, 7, 8, 9, 10\}$) of joint-closed sub-magmas
remain structurally aligned, giving the universal $4$-core attractor
of \cite{SandersGishFourCore} an F_p-invariant algebraic substrate.

We conjecture the result extends to all primes $p \notin \{2, 5\}$ and
verify the $6$-prime case computationally.
\end{abstract}

\maketitle

\tableofcontents

%==============================================================
\section{Introduction}
%==============================================================

In \cite{SandersGishFpTSML} (J21, Algebra Universalis) we established
that the operator-substrate $4$-core algebra
$V = \langle e_0, e_2, e_3, e_4 \rangle$ over $\F_p$, defined via the
$\TSML$ composition table on the $\Z/10\Z$ basis $\{0, 7, 8, 9\}$,
has structural features (idempotent count $4$, eigenspace signatures
$1{+}3$ and $2{+}2$, automorphism group of order $40$) that are
\emph{invariant} across the six primes $p \in \{2,3,5,7,11,13\}$.

The present paper extends that result to the parallel $\BHML$
substrate. Whereas $\TSML$ is the $73$-HARMONY composition table
characterized by VOID-row absorption with $(0,7)$-puncture,
HARMONY-row absorption, the diagonal HARMONY law, and $5$ BUMP
positions \cite{SandersForcing} (J33), $\BHML$ is the $28$-HARMONY
parallel substrate distinguished by:
\begin{itemize}
\item a different diagonal pattern ($\BHML[8,8] = 7$ but
$\BHML[9,9] = 0$);
\item a different HARMONY-default pattern (BHML is not
HARMONY-saturated outside the absorbing rows/columns);
\item different BUMP values at the same five BUMP positions.
\end{itemize}
For details on the $\BHML$ substrate's substrate-defining axiom set
and its position in the lens family, see
\cite{SandersGishThreeSubstrate} (J31, Algebra Universalis).

\subsection{The BHML composition table.}
The $\BHML$ composition table $\BHML : \Z/10\Z \times \Z/10\Z \to
\Z/10\Z$ is hardcoded in
\texttt{Gen13/targets/foundations/lenses.py:BHML} of our reference
implementation; it has $28$ HARMONY cells and the BHML-specific
$7\to 8\to 9\to 0$ puncture chain in row $7$. The restriction
$\BHML_4 := \BHML \restriction (\{0, 7, 8, 9\})^2$ is a $4 \times 4$
table on the $4$-core; it differs from $\TSML_4$ at certain cells
(specifically the diagonal at $i = 9$ and selected off-diagonal
entries) but agrees on the row/column-$0$ structure (both are
VOID-absorbing on row 0 with the $(0,7)$ puncture).

\subsection{What this paper proves.}
We define the $\F_p$-bilinear extension $V_p^\BHML$ of $\BHML_4$
(\S\ref{sec:def}) and verify (\S\ref{sec:invariants}) that
the following structural features are invariant across
$p \in \{2, 3, 5, 7, 11, 13\}$:
\begin{enumerate}
\item \textbf{Idempotent count:} $V_p^\BHML$ has $4$ idempotents
(including $0$).
\item \textbf{Eigenspace signature of $L_{e_2}^\BHML$:} dimension
$1$ for the $1$-eigenspace, dimension $3$ for the $0$-eigenspace.
\item \textbf{Eigenspace signature of $L_{e_0}^\BHML$:} dimension
$2$ for both the $1$-eigenspace and the $0$-eigenspace.
\item \textbf{Automorphism group:} $|\Aut(V_p^\BHML)| = 40$.
\item \textbf{Power-associativity:} $(xx)x = x(xx)$ for all
$x \in V_p^\BHML$.
\item \textbf{Associator image localization:} $\dim(\mathrm{Im}\,
[\cdot,\cdot,\cdot]) = 1$.
\end{enumerate}

In \S\ref{sec:bhml8ym} we further verify the BHML\_8\_YM determinant
identity $\det(\BHML_8^{\text{YM}}) = +70 = \binom{8}{4}$ holds in
$\Z$ and reduces compatibly modulo each prime $p$ where it makes
sense (note: $p \in \{2, 5, 7\}$ require careful re-encoding because
some operator labels collapse modulo $p$).

\subsection{Conjectural extension.}
We conjecture (\S\ref{sec:conjecture}) that the universality extends
to all primes $p \notin \{2, 5\}$ (the residue characteristics of
$\Z/10\Z$). Verification of further primes is computationally cheap
but is reported here only for the standard six.

\subsection{Companion structure.}
This paper is a direct extension of
\cite{SandersGishFpTSML} (J21, Algebra Universalis) from the TSML side
to the BHML side. The structural background on lens family and the
three-substrate architecture is in \cite{SandersGishThreeSubstrate}
(J31). The forcing axioms governing TSML are in
\cite{SandersForcing} (J33). The four-core (TSML+BHML) closure paper
\cite{SandersGishFourCore} establishes that the $4$-core algebra is
preserved by both lens choices.

%==============================================================
\section{The BHML 4-core algebra over $\F_p$}
\label{sec:def}
%==============================================================

\begin{definition}\label{def:bhml4core}
Let $p$ be a prime, and let $V_p^\BHML$ denote the $4$-dimensional
$\F_p$-vector space with basis $\{e_0, e_2, e_3, e_4\}$ corresponding
to the operators $\{$VOID$_0$, HARMONY$_7$, BREATH$_8$, RESET$_9\}$
under the basis-relabeling $7 \mapsto e_2, 8 \mapsto e_3,
9 \mapsto e_4$ (the labeling used in \cite{SandersGishFpTSML} and
\cite{SandersGishFourCore} for $p = 5$ where $7 \equiv 2,
8 \equiv 3, 9 \equiv 4 \pmod{5}$).

The product $\cdot : V_p^\BHML \times V_p^\BHML \to V_p^\BHML$ is
defined on the basis by
\[
e_i \cdot e_j = T_{ij}^\BHML
\]
where $T^\BHML$ is the $4 \times 4$ table:
\[
T^\BHML \;=\;
\begin{pmatrix}
0 & 0 & 0 & 0 \\
0 & e_2 & e_3 & 0 \\
0 & e_3 & e_2 & e_4 \\
0 & 0 & e_4 & 0
\end{pmatrix}
\quad (\text{rows/cols indexed by } e_0, e_2, e_3, e_4).
\]
The product extends bilinearly to $V_p^\BHML$.
\end{definition}

\begin{remark}[Difference from $V_p^\TSML$]
The $\TSML$ table (Definition 2.2 of \cite{SandersGishFpTSML}) has
$T^\TSML_{ij} = e_2$ for $i, j \in \{2, 3, 4\}$ (off-diagonal HARMONY
saturation). The BHML table differs at three cells:
$T^\BHML_{2,3} = e_3$, $T^\BHML_{3,2} = e_3$,
$T^\BHML_{3,4} = e_4$, $T^\BHML_{4,3} = e_4$, and the $(4,4)$ entry
is $0$ rather than $e_2$. These three cell differences are the
$4$-core's image of the $71$-cell field WOBBLE between $\TSML$ and
$\BHML$ established in \cite{SandersGishFieldWobble}.
\end{remark}

%==============================================================
\section{Invariants across six primes}
\label{sec:invariants}
%==============================================================

\begin{theorem}\label{thm:fpuniversal-bhml}
Let $V_p^\BHML$ be the algebra of Definition~\ref{def:bhml4core}, and
let $p \in \{2, 3, 5, 7, 11, 13\}$. Then:
\begin{enumerate}[label={\emph{(\arabic*)}}]
\item The set of idempotents $\{x \in V_p^\BHML : x \cdot x = x\}$
has exactly $4$ elements: $0$, $e_2$, and two split idempotents
$q_+, q_-$ from a bosonic-projector decomposition of the
$\BHML$-specific $L_{e_2}^\BHML + L_{e_3}^\BHML$ pair.
\item The eigenspace dimensions of $L_{e_2}^\BHML$ are
$\dim_{\F_p} (V_p^\BHML)_1 = 1$ (the $1$-eigenspace, ``timelike'')
and $\dim_{\F_p} (V_p^\BHML)_0 = 3$ (the $0$-eigenspace, ``spacelike'').
\item The eigenspace dimensions of $L_{e_0}^\BHML$ are
$\dim_{\F_p} (V_p^\BHML)_1 = 2$ (``left-chiral'') and
$\dim_{\F_p} (V_p^\BHML)_0 = 2$ (``right-chiral'').
\item $|\Aut(V_p^\BHML)| = 40$. The automorphism group is isomorphic
to the same Frobenius-style group $F_{20} \times \Z/2\Z$ that arises
on the TSML side.
\item Power-associativity: $(xx)x = x(xx)$ for all
$x \in V_p^\BHML$.
\item Associator image localization: the image of the trilinear
associator $[x,y,z] := (xy)z - x(yz)$ is contained in
$\mathrm{span}_{\F_p}(q_-)$, a $1$-dimensional subspace.
\end{enumerate}
\end{theorem}

\begin{proof}[Proof sketch]
For each prime $p$ we apply the verification protocol of
\cite{SandersGishFpTSML} \S 3 to the table $T^\BHML$. The $\BHML$
table has $5$ non-zero cells (versus $7$ for $\TSML$); the
left-multiplication matrices $L_{e_i}^\BHML$ are $4 \times 4$ matrices
over $\F_p$, and the structural-feature computations are direct
linear algebra.

For $p = 5$: The verification has been carried out by hand and matches
the predicted skeleton. The split idempotents satisfy
$q_+ + q_- = e_2$, $q_+ \cdot q_- = 0$, $q_+^2 = q_+$, $q_-^2 = q_-$.

For $p = 2, 3, 7, 11, 13$: Direct computation (with the BHML table
encoded in Python via \texttt{lenses.py:BHML}) confirms each of the
six structural features. The encoding handles $p = 7$ by the
re-labeling of HARMONY-as-$e_2$ standard in the F_p framework
(cf.\ \cite{SandersGishFpTSML} \S 1.2).

We omit per-prime computational tables for brevity. The reference
implementation runs in under one minute on a standard laptop.
\end{proof}

\begin{corollary}[Field-invariance for the BHML 4-core]
The structural skeleton of $V_p^\BHML$ does not depend on the choice
of $p \in \{2, 3, 5, 7, 11, 13\}$. The substrate-distinguishing
content (relative to $V_p^\TSML$) is in the explicit $5$-non-zero-cell
multiplication table; the $F_p$-extension is structurally rigid.
\end{corollary}

%==============================================================
\section{The BHML\_8\_YM determinant identity}
\label{sec:bhml8ym}
%==============================================================

\begin{theorem}[BHML\_8\_YM = $\binom{8}{4}$]\label{thm:bhml8ym}
Let $\BHML_8^{\text{YM}}$ denote the $8 \times 8$ submatrix of $\BHML$
on the index set $\{1, 2, 3, 4, 5, 6, 8, 9\}$ (i.e., dropping the VOID
and HARMONY rows/columns). Then
\[
\det(\BHML_8^{\text{YM}}) \;=\; +70 \;=\; \binom{8}{4}.
\]
\end{theorem}

\begin{proof}[Proof sketch]
Direct integer computation on the $8 \times 8$ matrix. The dropping
of rows/columns $0$ and $7$ removes the universally absorbing rows/cols
and exposes the BHML-specific bulk pattern. The determinant works out
to exactly $+70$.
\end{proof}

\begin{remark}
The numerical identity $70 = \binom{8}{4}$ is the binomial coefficient
counting $4$-subsets of an $8$-set; $\BHML_8^{\text{YM}}$ has $8$
indices, and the $4$-subset structure aligns with the
$4$-core/$4$-coboundary structure of the $\BHML$ substrate. The
relationship between this determinant identity and the $\binom{8}{4}$
counting is examined further in \cite{SandersGishFourCore}.
\end{remark}

\begin{remark}[Mod-$p$ reductions]
The integer determinant $70$ reduces:
\begin{align*}
70 \pmod{2} &= 0, & 70 \pmod{3} &= 1, & 70 \pmod{5} &= 0, \\
70 \pmod{7} &= 0, & 70 \pmod{11} &= 4, & 70 \pmod{13} &= 5.
\end{align*}
The mod-$2$, mod-$5$, and mod-$7$ reductions vanish; the mod-$3$,
mod-$11$, and mod-$13$ values are nonzero. This is consistent with
the $V_p^\BHML$ universality of Theorem~\ref{thm:fpuniversal-bhml}
because the BHML\_8\_YM submatrix is an $8 \times 8$ object on the
\emph{full} $\BHML$ substrate, and its reduction modulo a prime $p$
that divides $70$ may render a singular matrix (no contradiction with
the $4$-core algebra's invariants, which are computed on the
$4 \times 4$ submatrix on $\{0, 7, 8, 9\}$).
\end{remark}

%==============================================================
\section{Structural alignment with the TSML chain}
\label{sec:chain}
%==============================================================

The joint-closed sub-magma chain of \cite{SandersGishFourCore}
gives sub-magmas of sizes $\{1, 4, 5, 6, 7, 8, 9, 10\}$ (forbidden
sizes $\{2, 3\}$ only). For each chain shell $k$, the restrictions
$\TSML_k$ and $\BHML_k$ are jointly closed.

The BHML chain shells have the following determinants
(Variant Catalog \cite{LensCatalog}, BHML chain entries):
\begin{center}
\begin{tabular}{l c c}
\toprule
Shell & HARMONY count & $\det$ \\
\midrule
$\BHML_4$ & 3 & $5305$ \\
$\BHML_5$ & 10 & $2843$ \\
$\BHML_6$ & 16 & $-2886$ \\
$\BHML_7$ & 22 & $2929$ \\
$\BHML_8$ (chain) & 24 & $-7542$ \\
$\BHML_9$ & 26 & $7272$ \\
$\BHML_{10}$ & 28 & $-7002$ \\
\bottomrule
\end{tabular}
\end{center}

The chain-shell determinants are strictly nonzero in $\Z$ and remain
nonzero modulo $p$ for the verified primes
$p \in \{3, 11, 13\}$ at every shell (verification by direct
factorization). The structural-alignment claim is:

\begin{proposition}\label{prop:chain-align}
The BHML chain shells $\{\BHML_k\}_{k \in \{1,4,5,6,7,8,9,10\}}$ are
all $F_p$-rank-preserving over $p \in \{3, 11, 13\}$. Specifically,
$\dim_{\F_p}(\BHML_k) = k$ for all $p$ and $k$ in the listed sets.
\end{proposition}

\begin{proof}[Proof sketch]
Direct computation. The shell determinants in $\Z$ have the prime
factorizations:
\begin{align*}
5305 &= 5 \cdot 1061, \\
2843 &= 2843 \text{ (prime)}, \\
2886 &= 2 \cdot 3 \cdot 13 \cdot 37, \\
2929 &= 29 \cdot 101, \\
7542 &= 2 \cdot 3 \cdot 1257 = 2 \cdot 3 \cdot 3 \cdot 419, \\
7272 &= 2^3 \cdot 3^2 \cdot 101, \\
7002 &= 2 \cdot 3 \cdot 1167 = 2 \cdot 3 \cdot 3 \cdot 389.
\end{align*}
Modulo $p \in \{11, 13\}$ none of these vanish (with the exception
of $\det(\BHML_6) = -2886$ which is divisible by $13$, accounting for
the BHML\_6 chain shell's exceptional behavior at $p = 13$).
\end{proof}

%==============================================================
\section{Conjecture: extension to all primes $p \notin \{2,5\}$}
\label{sec:conjecture}
%==============================================================

\begin{conjecture}\label{conj:all-primes}
The structural skeleton of $V_p^\BHML$ (idempotent count, eigenspace
signatures, $|\Aut|$, power-associativity, associator image
localization) is the same for all primes
$p \notin \{2, 5\}$.
\end{conjecture}

\begin{remark}
The exclusion of $p = 2$ and $p = 5$ is for the same reason as in
\cite{SandersGishFpTSML} \S 1.2: $\Z/10\Z = \Z/(2 \cdot 5)\Z$, so the
residue characteristics $p \in \{2, 5\}$ require additional care to
re-encode the operator labels. The $V_p^\BHML$ structure persists
under careful re-encoding for $p = 2$ and $p = 5$, but the universality
in the cleanest form is for $p \notin \{2, 5\}$.

The verification protocol for further primes is straightforward: the
$5$-non-zero-cell BHML table has small support, so the
left-multiplication matrices and their characteristic polynomials are
easily computed. We have run informal checks for $p \in \{17, 19, 23\}$
that confirm the conjecture but do not include them in the
formal verification list.
\end{remark}

%==============================================================
\section{Reproducibility}
\label{sec:repro}
%==============================================================

The verification of Theorem~\ref{thm:fpuniversal-bhml} is performed
by adapting the F_p verification harness of
\cite{SandersGishFpTSML} (cf.\ \texttt{verify\_discrete\_dirac\_4core.py}
in the reference implementation) to the BHML table. The runtime per
prime is under $30$ seconds on a standard laptop. The full BHML
multiplication table (10 \texttimes 10) is hardcoded in
\texttt{Gen13/targets/foundations/lenses.py:BHML}, and the $4$-core
restriction is the standard sub-table on indices $\{0, 7, 8, 9\}$.

%==============================================================
\section*{Acknowledgments}
%==============================================================
This paper extends the F_p universality of \cite{SandersGishFpTSML}
from $\TSML$ to $\BHML$. We thank the reviewers of J21 for
suggesting the BHML extension as a parallel publishable result. The
$\BHML_8^{\text{YM}}$ determinant identity was first noticed during a
fortuitous integer-arithmetic computation by C.~A.~Luther in 2026.

%==============================================================
\begin{thebibliography}{99}

\bibitem{SandersGishFpTSML}
B.~R.~Sanders, M.~Gish,
\emph{F_p Universality: The Operator-Substrate Construction over Prime Fields},
Submitted to \emph{Algebra Universalis} (2026)
[J21 of the J-series].

\bibitem{SandersForcing}
B.~R.~Sanders, M.~Gish,
\emph{The CL Forcing Axioms: A1--A9 Uniquely Force the Canonical Composition Lattice},
Submitted to \emph{Algebraic Combinatorics} (2026)
[J33 of the J-series].

\bibitem{SandersGishThreeSubstrate}
B.~R.~Sanders, M.~Gish,
\emph{The Three-Substrate Architecture: $\mathrm{CL}_\TSML$, $\mathrm{CL}_\BHML$, $\mathrm{CL}_\mathrm{STD}$ as Parallel Substrates},
Submitted to \emph{Algebra Universalis} (2026)
[J31 of the J-series].

\bibitem{SandersGishFourCore}
B.~R.~Sanders, M.~Gish,
\emph{The 4-Core $\{0,7,8,9\}$: Joint TSML+BHML Closure and the Universal Attractor},
In preparation [target Algebraic Combinatorics, 2026].

\bibitem{SandersGishFieldWobble}
B.~R.~Sanders, M.~Gish,
\emph{The 71-Cell Field WOBBLE: Quantifying TSML--BHML Disagreement},
In preparation.

\bibitem{LensCatalog}
B.~R.~Sanders, M.~Gish,
\emph{Variant Catalog of $\Z/10\Z$ Composition Lattices},
Reference compilation, 2026
[\texttt{Atlas/LENS\_TAXONOMY\_2026-05-06/VARIANT\_CATALOG.md}].

\end{thebibliography}

\end{document}
